\section{Cyclic covers and realized symmetries}\label{sec:cyclic}

Two criteria organize the passage from a filling obstruction to the
homeomorphism conclusions.  The first forbids a prescribed outer
automorphism from being induced by a homeomorphism.  The second uses
finite symmetry groups to distinguish two manifold models without
prescribing the homotopy class of a possible homeomorphism.

\subsection{A criterion from cyclic covers}

We first isolate the passage from manifold existence in finite covers to
a homotopy equivalence that cannot be represented by a homeomorphism.
This passage is independent of the four-dimensional construction.
A \emph{closed topological $n$-manifold model} of a space is a closed,
connected topological $n$-manifold equipped with a homotopy equivalence
to that space.

\begin{theorem}[Cyclic-cover criterion]\label{cyc:criterion}
Let $Q$ be a connected finite aspherical CW complex, put
$G=\pi_1(Q)$, and let $w\colon G\twoheadrightarrow\mathbb Z$ be an
epimorphism.  For every positive integer $d$, write
\[
 G_d=w^{-1}(d\mathbb Z)
 \qquad\text{and}\qquad
 Q_d\longrightarrow Q
\]
for the associated connected cyclic cover.  Fix an integer $n\geq1$.
Suppose that $Q_2$ has a closed topological $n$-manifold model $M$, but
that $Q_d$ has no such model for every sufficiently large odd integer
$d$.

Choose $\gamma\in G$ with $w(\gamma)=1$, and use the model equivalence
to identify $\pi_1(M)$ with $G_2$.  The outer automorphism of $G_2$
represented by
\[
 \alpha(g)=\gamma g\gamma^{-1}
\]
is induced by a self-homotopy equivalence of $M$.  No homeomorphism of
$M$ induces this outer automorphism.  In particular, that
self-homotopy equivalence is not homotopic to a homeomorphism.
\end{theorem}

\begin{proof}
The manifold $M$ is aspherical, since it is homotopy equivalent to
$Q_2$.  Topological manifolds have CW homotopy type~\cite{Milnor1959}.
The classification
of maps between Eilenberg--Mac Lane spaces therefore realizes $\alpha$
by a map $f\colon M\to M$, and realizes its inverse by a homotopy
inverse to $f$; see \cite[Proposition~1B.9]{Hatcher2002}.

Suppose that a homeomorphism $h\colon M\to M$ induces the outer class
of $\alpha$.  We will construct closed manifold models for all
sufficiently large odd-degree covers of $Q$.
Put $K=\ker w$, and let
\[
 X=K\backslash\widetilde M
\]
be the infinite cyclic cover of $M$, where $\widetilde M$ is its
universal cover.  We use left deck actions throughout.  The element
$\gamma^2\in G_2$ induces a deck transformation $T$ of $X$, and $T$
generates its deck group over $M$.

Choose a lift $u\colon\widetilde M\to\widetilde M$ of $h$ that
intertwines the deck action by exactly $\alpha$:
\begin{equation}\label{cyc:exact-lift}
 u g u^{-1}=\gamma g\gamma^{-1}
 \qquad(g\in G_2).
\end{equation}
Indeed, any lift induces an automorphism in the prescribed outer
class.  Composing that lift with a deck transformation removes the
inner discrepancy.  Since $\alpha(K)=K$, the map $u$ descends to a
homeomorphism $h'\colon X\to X$.  Moreover,
$\alpha(\gamma^2)=\gamma^2$, so \eqref{cyc:exact-lift} gives
\begin{equation}\label{cyc:commutation}
 h'T=Th'.
\end{equation}

We next show that a sufficiently large translation composed with $h'$
acts freely and cocompactly on $X$.  Represent the homomorphism
$w/2\colon G_2\to\mathbb Z$ by a map $M\to S^1$ and lift that map
to obtain a continuous height function
\[
 \tau\colon X\longrightarrow\mathbb R,
 \qquad \tau(Tx)=\tau(x)+1.
\]
This height is proper.  To see this, choose a compact set $C\subset X$
whose $T$-translates cover $X$, using compactness of $M$, and choose
$a,b\in\mathbb R$ with $\tau(C)\subset[a,b]$.  A bounded closed
height band $\tau^{-1}([A,B])$ is a closed subset of
\[
 \bigcup_{\substack{j\in\mathbb Z\\ A-b\leq j\leq B-a}}T^jC,
\]
which is a finite union of compact sets.

By \eqref{cyc:commutation}, the continuous function
$x\mapsto\tau(h'x)-\tau(x)$ is $T$-invariant.  It is therefore
bounded, say in absolute value by $c\geq0$.  For an integer $r>c+1$,
put
\[
 F_r=T^rh'.
\]
For every $x\in X$ we have
\begin{equation}\label{cyc:height-increment}
 0<r-c\leq\tau(F_rx)-\tau(x)\leq r+c.
\end{equation}
Consequently no nonzero power of $F_r$ fixes a point.  Iterating the
lower bound shows that only finitely many translates of any compact
set can meet that set; thus the $\mathbb Z$-action generated by $F_r$
is properly discontinuous.  Finally, along every orbit the height
tends to $+\infty$ in positive time and to $-\infty$ in negative
time.  The first nonnegative height on such an orbit belongs to
$[0,r+c]$, by the upper bound in \eqref{cyc:height-increment}.
The compact band $\tau^{-1}([0,r+c])$ therefore meets every orbit.
This proves cocompactness.

It follows that
\[
 Y_r=X/\langle F_r\rangle
\]
is a closed connected topological $n$-manifold.  Its universal cover
is $\widetilde M$, so $Y_r$ is aspherical.  To identify its fundamental
group, lift $F_r$ to
\[
 \widetilde F_r=\gamma^{2r}u
 \quad\text{on }\widetilde M.
\]
Equation \eqref{cyc:exact-lift} gives
\begin{equation}\label{cyc:odd-action}
 \widetilde F_r k\widetilde F_r^{-1}
   =\gamma^{2r+1}k\gamma^{-(2r+1)}
 \qquad(k\in K).
\end{equation}
The deck group over $Y_r$ is generated by $K$ and $\widetilde F_r$.
Its quotient by $K$ is infinite cyclic, since its action on $X$ is
generated freely by $F_r$.  Hence \eqref{cyc:odd-action} identifies it
with
\[
 \pi_1(Y_r)
   \cong K\rtimes_{\operatorname{Ad}_{\gamma^{2r+1}}}\mathbb Z
   \cong G_{2r+1}.
\]
For the second isomorphism, send the generator of $\mathbb Z$ to
$\gamma^{2r+1}$ and act identically on $K$.  This is bijective because
every element of $G_{2r+1}$ has a unique expression
$k\gamma^{(2r+1)j}$ with $k\in K$ and $j\in\mathbb Z$.

Both $Y_r$ and $Q_{2r+1}$ are Eilenberg--Mac Lane spaces, so this
group isomorphism is realized by a homotopy equivalence.  Taking $r$
sufficiently large contradicts the assumed absence of odd-degree
manifold models.
\end{proof}

\begin{remark}\label{cyc:homotopy-involution}
The equivalence $f$ in Theorem~\ref{cyc:criterion} satisfies
$f^2\simeq\operatorname{id}_M$.  Indeed,
$\alpha^2=\operatorname{Ad}_{\gamma^2}$ is inner because
$\gamma^2\in G_2$, and unbased
homotopy classes of maps between aspherical CW-type spaces are
classified by homomorphisms modulo inner automorphisms of the target.
The proof nevertheless imposes no order condition on a hypothetical
realizing homeomorphism.  Its only geometric input is the bounded
height displacement of a lift commuting with the deck translation.
\end{remark}

\subsection{A criterion for distinct manifold models}

Let $X$ be a connected topological manifold and choose an isomorphism
$\mu_X:\pi_1(X)\xrightarrow{\cong}\Pi$ to a group $\Pi$.  Basepoint
paths affect an induced automorphism only by an inner automorphism, so
this marking defines a subgroup
\[
 R_X=\operatorname{im}\bigl(\operatorname{Homeo}(X)
                  \longrightarrow\Out(\Pi)\bigr).
\]
Changing the marking conjugates $R_X$ in $\Out(\Pi)$.  More generally,
if $X$ and $Y$ are so marked and $f:X\to Y$ is a homeomorphism, the
outer isomorphism $\rho$ defined by $\mu_Y f_*\mu_X^{-1}$ satisfies
\begin{equation}\label{sym:transport}
 R_Y=\rho R_X\rho^{-1}.
\end{equation}
This follows by transporting self-homeomorphisms through $f$.

\begin{proposition}[A Sylow criterion for manifold models]
\label{sym:criterion}
Let $X$ and $Y$ be connected topological manifolds whose fundamental
groups are marked by a group $\Pi$, and suppose that $\Out(\Pi)$ is
finite.  Let $\mathsf A,\mathsf B\leq\Out(\Pi)$ be $2$-subgroups of
the same order.  Suppose that there is an involution $\delta$ for which
\[
 N_{\Out(\Pi)}(\mathsf A)=\mathsf A\times\langle\delta\rangle
\]
is an internal direct product, and that
\[
 \mathsf A\leq R_X,\qquad \mathsf B\leq R_Y,
 \qquad \delta\notin R_X.
\]
If $\mathsf A$ and $\mathsf B$ are not conjugate in $\Out(\Pi)$, then
$X$ and $Y$ are not homeomorphic.
\end{proposition}

\begin{proof}
Since $\mathsf A\leq R_X$, the inclusion of any $a\delta$ with
$a\in\mathsf A$ in $R_X$ would imply $\delta\in R_X$.  Hence
\[
 N_{R_X}(\mathsf A)
 =R_X\cap\bigl(\mathsf A\times\langle\delta\rangle\bigr)
 =\mathsf A.
\]
Choose a Sylow $2$-subgroup $P$ of $R_X$ containing $\mathsf A$.
A proper subgroup of a finite $2$-group is strictly contained in its
normalizer.  For completeness, this follows by induction on the group
order: an element of the center outside the subgroup enlarges its
normalizer, while if the center is contained in the subgroup the claim
follows from the induction hypothesis in the quotient by the center.
Thus $\mathsf A<P$ would contradict the displayed equality, and
$\mathsf A$ is a Sylow $2$-subgroup of $R_X$.

If a homeomorphism $X\to Y$ existed, \eqref{sym:transport} would place
a conjugate of $\mathsf B$ inside $R_X$.  Its order is $|\mathsf A|$,
so it too would be a Sylow $2$-subgroup of $R_X$.  Sylow conjugacy in
$R_X$ would make $\mathsf A$ and $\mathsf B$ conjugate in $\Out(\Pi)$,
contrary to the hypothesis.
\end{proof}
